\section{Introduction}
\IEEEPARstart{V}{isual-language} tracking (VLT) localizes a target object in a video sequence given both an initial visual template and a natural language description. Compared with conventional visual tracking, language provides complementary semantic cues that are especially useful under appearance changes, occlusion, distractors, and ambiguous scenes.

Despite recent progress, existing VLT methods still face a key limitation: the language branch is usually used as a conditioning signal, while the learned visual representation is not explicitly constrained to be semantically consistent with the textual description. As a result, the tracker may fit visual localization objectives without learning a text-aware representation that can be inspected or explained.

We propose SALTTrack, a semantic alignment and lightweight tuning framework for visual-language tracking. Our main contribution is a text-guided semantic alignment loss that supervises the relation between predicted visual features, ground-truth visual features, and target text features. Instead of simply maximizing absolute visual-text similarity, we emphasize direction-aware semantic correction, which encourages the prediction feature to move toward the semantic direction implied by the ground truth and the language description. A mild distance term is further used to improve semantic compactness without forcing collapse.

In addition, we introduce a routed LoRA adaptation mechanism for efficient fine-tuning. The routed design dynamically modulates low-rank adaptation paths according to semantic reliability, allowing the tracker to adapt to VLT data without updating the full model.

The main contributions are summarized as follows:
\begin{itemize}
    \item We introduce a semantic alignment objective for visual-language tracking, providing explicit text-guided supervision beyond box regression and classification losses.
    \item We formulate alignment as direction-dominant semantic correction with controlled distance regularization, improving interpretability while avoiding over-constrained text-visual similarity.
    \item We develop a semantic-guided routed LoRA tuning strategy for parameter-efficient adaptation.
    \item We conduct experiments on representative VLT benchmarks and provide ablations and visual analyses to validate both effectiveness and interpretability.
\end{itemize}
